\section{Gmapping SLAM library}

Gmapping is a open-source library implementing two-dimensional simultaneous localization and mapping (SLAM). Odometry data and range data from a laser scanner are used as input for the mapping. 

Gmapping uses a Rao-Blackwellized Particle Filter. SLAM is a difficult problem as for mapping the pose of the robot is needed and for accurate localization a map is needed. Rao-Blackwellized mapping systems use a particle filter to estimate the distribution over possible trajectories. Each particle holds its own map, which is created with occupancy grid mapping with known pose. The particle filter algorithm used consists of three steps, sampling, importance weighting and resampling. In the sampling step poses of the particles are updated by a movement step sampled from a proposal distribution. This proposal distribution is based on the odometry data and in more sophisticated systems further sensor data. In the importance weighting step a weight is assigned to every particle representing the likelyhood of the particle. In the resampling step particles with low weight are replaced by samples with a higher weight. This step is needed as a continuous distribution is approximated by only a finite number of particles. After the resampling the maps of the particles are updated.

The problem of this approach is the high computational cost when using a large number of particles, which is usually needed for a good representation of the distribution. Therefore Gmapping also takes the range measurements into account in the sampling step. This way unlikely particle poses can be ruled out immediately and therefore a much lower particle number is sufficient. To avoid the problem of particle depletion Gmapping also uses a sophisticated resampling technique.